\documentclass{ximera}

\begin{document}

    \title{Lecture 11}
    
    \begin{abstract}  
    Outline: \\
    - Cantor set \\
    - Open sets
\end{abstract}

\maketitle

Cantor set:

\begin{align*}
C_0 &= [0,1]    \\
C_1 &= [0,\frac13] \cup [\frac23,1]
\\
C_2 &= \left( [0,\frac19] \cup [\frac29,\frac13]\right) \cup \left( [\frac23,\frac79] \cup [\frac89,1]\right)\\
\vdots
\end{align*}
Cantor set: $\bigcap\limits_{n=0}^\infty C_n$

\begin{itemize}
\item Is the Cantor set empty?  \textcolor{red}{No, the NIP tells us it is nonempty!}
\item How many elements do you think are inside of the Cantor set? \textcolor{red}{It turns out infinitely many, indeed uncountably many, as you can identify each set with the direction you go (left = 0, right = 1) based on the previous subinterval you start with, and thus you can get as many sequences of $1$'s and $0$'s as you like; if you want to study cardinality in more depth, go back to Section 1.5, and do some of the exercises there.}
\item What is the length of the Cantor set? \textcolor{red}{ Let's compute the length of the removed intervals:
\begin{align*}
    \frac13 + 2\cdot\frac19 + 4 \cdot \frac{1}{27} + \cdots + 2^{n-1}\frac{1}{3^n}
\end{align*}
This is a geometric series!  In this case, $a = \frac13$ and r = $\frac23$, and therefore the series converges to $\frac{a}{1-r} = \frac{1/3}{1-2/3} = 1$.  So, the Cantor set has $0$ length (1-1 =0).}
\end{itemize}

Behavior at infinity is strange.

The Cantor set is what is known as a \textit{fractal}.  It lives in a weird world, a world we can say is not quite 1-dimensional, but it still has so many elements in it!  Let's agree that a point has dimension $0$.  What happens when I magnify different structures by a factor of $3$? A point is still a point, it does not change length.  A line increases by a factor of $3$ in length.  A square increases by a factor of $3^2$ in area, and a cube increases by a factor $3^3$ in volume, etc etc.  This is our usual, while unrigorous, notion of dimension.  What happens to the Cantor set?  $3C_0 = [0,3]$; deleting the middle third gives $[0,1]\cup[2,3]$, which is going to give us two copies of the original Cantor set! One for $[0,1]$, and the other for $[2,3]$.  So, this implies that $2 = 3^x$, or $x =  \ln{2}/\ln{3} \approx 0.631$.  So the Cantor set, from this particular notion of dimension (of which there are many), is not fully 1D!!!

%\textcolor{red}{Pull up an article on fractals in real life.}

Notions of behavior at infinity are, as we have noticed in this course, NOT straightforward, and give us WONDERFUL ways to perceive and understand reality!  The general concept of infinity is baked into what we have at this point very informally called ``open" and ``closed" sets. Let's define them now.

Let $a \in \mathbb{R}$ and $\epsilon > 0$.  Recall that an $\epsilon$-neighborhood of $a$ is the set
\begin{align*}
    V_{\epsilon}(a) = \{x \in \mathbb{R}:|x-a|<\epsilon\}.
\end{align*}

\begin{definition}
    A set $O\subset\mathbb{R}$ is open if for all points $a \in O$, there exists an $\epsilon$- neighborhood $ V_{\epsilon}(a)  \subset{O}$.
\end{definition}

\begin{exercise}
    Conjecture some open sets in $\mathbb{R}$, and conjecture some sets that cannot be open.
\end{exercise}

\begin{theorem}
    \begin{enumerate}
        \item The union of an arbitrary collection of open sets is open.
        \item The intersection of a finite collection of open sets is open.
    \end{enumerate}
\end{theorem}

\textcolor{red}{Why is the second assumption here that the intersection has to be finite?}

\begin{proof}
\begin{enumerate}
   \item  Let $\{O_\lambda:\lambda \in \Lambda\}$ be a collection of open sets and let $O = \bigcup\limits_{\lambda\in \Lambda} O_\lambda$. Let $a \in O.$  Since $a \in O$,, there exists $\lambda'$ such that $V_\epsilon(a) \subset O_{\lambda'}\subset O$.
\item Let $\{O_1, ..., O_N\}$ be a finite collection of open sets.  If $a \in \bigcap\limits_{n=1}^N O_n$, then $a \in O_i$ for all $i = 1,..., n$.  For each set, there exists an $\epsilon_i$ such that $V_{\epsilon_i}(a) \subset O_i$. 
%\textcolor{red}{Draw pictures here.} 
Then let $\epsilon = \min{i=1,..., n} \epsilon_i$.  Then, $V_\epsilon(a) \subset V_{\epsilon_i}(a)$ for all $i =1,...,n$, and therefore $V_\epsilon(a) \subset O_i$ for all $i=1,..., n$.
\end{enumerate}
\end{proof}


\end{document}